% ===== tr-04 块2：第 4 章 pp.14--26；本片 pp.19--22（原文页 19--22）=====
% [原文 p.19]
为了使 SYZ 猜想的全强度版本成立，我们在第 3 节所见的拓扑环面纤维化的强对偶形式应当在特殊拉格朗日层面继续成立。这将意味着一个镜像对 $X,\check X$ 会具有特殊拉格朗日环面纤维化 $f:X\to B$ 与 $\check f:\check X\to B$，其判别轨迹为余维二，而且 $f$ 与 $\check f$ 的判别轨迹会重合。这些纤维化随后在判别轨迹之外互为对偶。

在非紧环面簇 $X$ 上存在一些特殊拉格朗日纤维化的例子，其判别轨迹看起来与我们已在拓扑情形中描述过的非常相似。特别地，若 $X$ 是一个 $n$ 维 Ricci 平坦 Kähler 流形，带有一个保持度量与全纯 $n$-形式的 $T^{n-1}$-作用，则 $X$ 会有一个非常好的特殊拉格朗日纤维化，其判别轨迹余维二。（见 [21] 与 [16]。）然而，Dominic Joyce（见 [48] 及其中引用的其他论文）开始研究一些三维的 $S^1$-不变例子，发现了相当不同的行为。[19] 中有一个论证表明，若特殊拉格朗日纤维化是 $C^\infty$ 的，则判别轨迹将是（Hausdorff）余维二的。然而，Joyce 发现了不可微、只是分段可微的例子，而且其判别轨迹是余维一的：

\begin{example}{4.5.1}
定义 $F:\C^3\to\R\times\C$ 为 $F(z_1,z_2,z_3)=(a,c)$，其中 $2a=|z_1|^2-|z_2|^2$，且
\[
c=\begin{cases}
z_3 & a=z_1=z_2=0\\
z_3-\bar z_1\bar z_2/|z_1| & a\ge0,\ z_1\ne0\\
z_3-\bar z_1\bar z_2/|z_2| & a<0.
\end{cases}
\]
\end{example}
% [Tutor Note: 分段定义的两个分母分别为 |z_1|（a≥0 且 z_1≠0 支）与 |z_2|（a<0 支），恰好对应把 R²×S¹ 中沿哪一族圆坍缩；已在 700 dpi 下逐一核对。另注意 2a=|z_1|²−|z_2|² 中的系数 2 与第一行的 a=z_1=z_2=0 不可省。]

容易看出，若 $a\ne0$，则 $F^{-1}(a,c)$ 同胚于 $\R^2\times S^1$；而若 $a=0$，则 $F^{-1}(a,c)$ 是 $T^2$ 上的一个锥：本质上，$\R^2\times S^1$ 中的一个 $S^1$ 坍缩成一点。此外，这个映射的所有纤维都是特殊拉格朗日的，而它显然只是分段光滑的。判别轨迹就是由 $a=0$ 给出的整个平面。

这个例子迫使人们重新评估 SYZ 猜想的强形式。在进一步的工作中，Joyce 找到了证据，支持三维中一般特殊拉格朗日纤维化的一个更可能的图景。判别轨迹将不再是余维二的图，而是余维一的一块（blob）。通常，奇异纤维的奇点之并会是一个 Riemann 曲面，而它会映射到 $B$ 中一个 amoeba 形状的集合，即判别轨迹看起来像图 4.4 中右边而不是左边的图，并且会是旧的余维二判别轨迹图景的一个加厚。

Joyce 还作了一些论证，提示这个加厚的判别轨迹在第 3 节所见的
% [原文 p.20]
\figimg[0.75]{p4_fig04.png}
\figtitle{4.4}{}
% [Tutor Note: 原文图注仅有编号 "Figure 4."，无描述文字，故图题第二参数留空。左图为旧的余维二判别轨迹（三价图的奇点）、右图为 Joyce 预期的余维一加厚；正文说「像右边而不是左边」。]
两类基本顶点的邻域中必然看起来根本不同，而两类顶点预期大致如图 4.4 所示。于是上面提到的强对偶形式------即我们预期镜像对上特殊拉格朗日纤维化的判别轨迹相同------不可能成立。若情形如此，人们就需要用较弱的形式取代这一强对偶形式。

看起来，重述 SYZ 猜想的最好方式很可能是取极限形式。我们目前所理解的镜像对称与 Calabi--Yau 流形的退化有关。给定圆盘 $D$ 上的平坦族 $f:\mathcal{X}\to D$，其 $0$ 上的纤维 $\mathcal{X}_0$ 奇异，而所有其他纤维是 $n$ 维 Calabi--Yau 流形，我们称该族为\emph{极大幂零}的，如果单值变换 $T:H^n(\mathcal{X}_t,\Q)\to H^n(\mathcal{X}_t,\Q)$（$t\in D$ 非零）满足 $(T-I)^{n+1}=0$ 但 $(T-I)^n\ne0$。镜像对称的一个标准预期是：镜像应当与 Calabi--Yau 流形的极大幂零退化相联系。特别地，给定某个 Calabi--Yau 流形的同一个复模空间中的两个不同的极大幂零退化，人们可能得到不同的镜像流形。这类退化在物理学文献中通常称为「大复结构极限」，尽管有时这个说法被用来对该退化附加一些额外条件，见 [63]。
% [Tutor Note: 极大幂零的定义是 (T−I)^{n+1}=0 但 (T−I)^n≠0（n = 纤维的复维数），指数与不等号方向不可写反；此定义用上同调 H^n，而 p22 中 W_0 的取像用同调 H_n，原文两处记号不同，照录。]

我们回忆 Gromov--Hausdorff 收敛的定义，这是一种度量空间序列的收敛概念。

\begin{definition}{4.5.2}
设 $(X,d_X)$、$(Y,d_Y)$ 是两个紧度量空间。假设存在映射 $f:X\to Y$ 与 $g:Y\to X$（不必连续），使得对所有 $x_1,x_2\in X$，
\[
|d_X(x_1,x_2)-d_Y(f(x_1),f(x_2))|<\epsilon,
\]
并且对所有 $x\in X$，
\[
d_X(x,g\circ f(x))<\epsilon,
\]
% [原文 p.21]
以及对 $Y$ 成立的两条对称性质也成立。则我们说 $X$ 与 $Y$ 之间的 Gromov--Hausdorff 距离至多为 $\epsilon$。Gromov--Hausdorff 距离 $d_{GH}(X,Y)$ 是所有这些 $\epsilon$ 的下确界。
\end{definition}
% [Tutor Note: 定义只显式写出了 X 方向的畸变条件（d_X 与 d_Y∘f 的比较、g∘f 接近恒等），而「对 Y 的两条对称性质」原文未重复写出；完整定义需同时要求 |d_X−d_Y| 双向与 d_Y(y, f∘g(y))<ε。]

由 Gromov 的结果（例如见 [68] 第 281 页 Cor. 1.11）可知，直径 $\le C$ 的紧 Ricci 平坦流形所构成的空间关于 Gromov--Hausdorff 距离是预紧的，即任何这样的流形序列都有子序列关于 Gromov--Hausdorff 距离收敛到一个度量空间。这个度量空间可能相当糟糕；这完全超出了代数几何的疆域！然而，这引出了如下自然的问题。给定 Calabi--Yau 流形的一个极大幂零退化 $\mathcal{X}\to D$，取收敛到 $0$ 的序列 $t_i\in D$，并考虑序列 $(\mathcal{X}_{t_i},g_{t_i})$，其中 $g_{t_i}$ 是选定的 Ricci 平坦度量，使得 $\operatorname{Diam}(g_{t_i})$ 保持有界。$(\mathcal{X}_{t_i},g_{t_i})$ 的 Gromov--Hausdorff 极限，或某个收敛子列的极限，是什么？

\begin{example}{4.5.3}
考虑由 $t$ 参数化的一族退化椭圆曲线，由 $\C/(\Z+\Z\tau)$ 给出，其中 $1$ 与 $\tau=\frac{1}{2\pi i}\log t$ 是这些椭圆曲线的周期。若让 $t$ 沿正实轴趋近 $0$，则我们可以直接把它看作一族周期为 $1$ 与 $i\alpha$ 的椭圆曲线 $\mathcal{X}_\alpha$，其中 $\alpha\to\infty$。若在 $\mathcal{X}_\alpha$ 上取标准 Euclidean 度量 $g$，则 $\mathcal{X}_\alpha$ 的直径无界。为得到有界直径，我们把 $g$ 换成 $g/\alpha^2$；等价地，我们可以在 $\C$ 上保持 $g$ 不变，而把椭圆曲线的周期改为 $1/\alpha$ 与 $i$。于是很清楚，这样的椭圆曲线序列的 Gromov--Hausdorff 极限是一个圆周 $\R/\Z$。
\end{example}
% [Tutor Note: τ = (1/2πi)log t（不是 2πi log t）；极限圆周长来自把周期换成 1/α 与 i 并令 α→∞，此时 1/α 方向的圆坍缩，剩下 R/Z。]

这个简单的例子启发人们提出关于极大幂零退化的第一个猜想，它分别由我与 Wilson 一方面 [38]、以及 Kontsevich 与 Soibelman 另一方面 [53] 独立提出。

\begin{conjecture}{4.5.4}
设 $\mathcal{X}\to D$ 是单连通、具有完整 $SU(n)$ 和乐群的 Calabi--Yau 流形的一个极大幂零退化，$t_i\in D$ 且 $t_i\to0$，并设 $g_i$ 是 $\mathcal{X}_{t_i}$ 上标准化使直径固定为 $C$ 的 Ricci 平坦度量。则 $(\mathcal{X}_{t_i},g_i)$ 的某个收敛子列收敛到一个度量空间 $(X_\infty,d_\infty)$，其中 $X_\infty$ 同胚于 $S^n$。此外，$d_\infty$ 由 $X_\infty\setminus\Gamma$ 上的一个 Riemann 度量诱导，其中 $\Gamma\subseteq X_\infty$ 是一个余维二的集合。
\end{conjecture}

这里极限的拓扑取决于非奇异纤维 $\mathcal{X}_t$ 的性质；例如，若 $\mathcal{X}_t$ 换成超 Kähler 的，则我们会预期极限是一个射影空间。另外，即使在完整 $SU(n)$ 和乐群的情形，若 $\mathcal{X}_t$ 不是单连通的，我们也会预期出现诸如 $\Q$-同调球这样的极限。

猜想 5.4 直接受 SYZ 猜想启发。假设我们有特殊拉格朗日纤维化 $f_i:\mathcal{X}_{t_i}\to B_i$。则当逼近极大幂零退化时，
% [原文 p.22]
可以看出这些纤维化的纤维体积趋于零。这提示这些纤维发生坍缩，有望把底留作极限。

这个猜想由我与 Wilson 于 2000 年在 [38] 中对 K3 曲面证明。证明依赖于 K3 曲面的一些令人愉快的事实。首先，它们是超 Kähler 流形，而特殊拉格朗日环面纤维化在复结构作一次超 Kähler 旋转之后变成椭圆纤维化。由于在 K3 曲面上构造椭圆纤维化很容易，而且这样的纤维化确实来自极大幂零退化的数据，因此很容易得到特殊拉格朗日纤维化。一旦做到这一点，就需要对 Ricci 平坦度量在极限中的行为作细致的分析。这是通过构造 Ricci 平坦度量的良好逼近来完成的，其中用到复维数二维时特殊拉格朗日纤维化的奇异纤维附近这些度量的显式局部模型的存在性。

所用的技术大多在高维不可用。然而，更近得多的时候，我与 V. Tosatti 和 Y. Zhang 在 [36] 的工作中得到了超 Kähler 情形下较弱的坍缩结果，其中假设存在类似于 K3 情形中椭圆纤维化的阿贝尔簇纤维化。我们没有得到显式的逼近 Ricci 平坦度量，而是利用了 Tosatti 在 [76] 中的先验估计。

在一般的 Calabi--Yau 情形，朝着该猜想唯一的进展是 Zhang 在 [78] 中的工作：它证明了在具有有界单射半径与截面曲率的 Calabi--Yau 流形区域中存在特殊拉格朗日纤维化，并由特殊拉格朗日纤维化的存在性推出局部坍缩。

SYZ 对猜想 5.4 的启发也给出了该猜想的一个极限形式。试图证明 Calabi--Yau 流形上特殊拉格朗日纤维化的存在性有许多问题。甚至在极大幂零退化附近是否存在单个特殊拉格朗日环面都未知，但我们预期当我们接近极大幂零点时会更容易找到它们。此外，即使我们找到一个特殊拉格朗日环面，我们知道它在一个 $n$ 维族中移动，但我们不知道它的形变是否填满整个流形。另外，即使填满，也不能保证我们得到流形的一个叶状结构：附近的特殊拉格朗日子流形可能相交。（例子见 [60]。）所以我们转而只看特殊拉格朗日环面的模空间。

给定 $n$ 维 Calabi--Yau 流形的一个极大幂零退化，已知 $(T-I)^n:H_n(\mathcal{X}_t,\Q)\to H_n(\mathcal{X}_t,\Q)$ 的像是一维子空间 $W_0$。假设给定满足 $i\to\infty$ 时 $t_i\to0$ 的序列 $t_i$，且当 $t_i$ 充分接近零时，存在一个生成 $W_0$ 的特殊拉格朗日 $T^n$。这正是我们预期找到与极大幂零退化相联系的特殊拉格朗日纤维化的纤维之处。设 $B_{0,i}$ 是这个环面的形变模空间；$B_{0,i}$ 的每个点对应于 $\mathcal{X}_{t_i}$ 中一个光滑的特殊拉格朗日环面。这个流形于是带有
% [Tutor Note: 这里 (T−I)^n 作用在 H_n（同调）上，与 p20 极大幂零定义中的 H^n（上同调）不同；W_0 是「极限混合 Hodge 结构」的一维部分，后文用它来挑出应作为纤维的特殊拉格朗日环面。]
